% Segundo informe OII464. Formato institucional de la Escuela de Ingeniería
% Informática PUCV y propuesta de selección adaptativa de movimientos.
\documentclass[12pt,letterpaper]{report}
\usepackage[spanish,es-nodecimaldot]{babel}
\usepackage{pucv_inf_2024}
\usepackage{amssymb,booktabs,array,microtype,csquotes}
\usepackage[hidelinks]{hyperref}
\addbibresource{referencias.bib}
\renewcommand{\thesection}{\arabic{section}}
\renewcommand{\thesubsection}{\thesection.\arabic{subsection}}
\renewcommand{\thetable}{\arabic{table}}
\captionsetup[table]{name=Tabla}
\setlength{\headheight}{14pt}
\hypersetup{
 pdftitle={Programación de cirugías electivas mediante búsqueda tabú y selección adaptativa de movimientos},
 pdfauthor={José Basualto, Cristopher Dauros, Carlos Ruz},
 pdfsubject={Segunda entrega del proyecto OII464},
 pdfkeywords={programación de cirugías, búsqueda tabú, selección adaptativa de operadores, aprendizaje en línea}
}

\begin{document}
\hypersetup{pageanchor=false}
% Adaptación de la portada para asignaturas de la plantilla PUCV 2024.
\title{Programación de cirugías electivas\\[0cm]
mediante búsqueda tabú y selección\\[0.2cm]
adaptativa de movimientos}
\author{José Basualto, Cristopher Dauros y Carlos Ruz}
\date{2 de noviembre de 2026}

\begin{titlepage}
  \makeatletter
  \thispagestyle{empty}
  \begin{picture}(0,0)
    \put(-75,-80){\includegraphics[width=20cm]{Portadas/imagenes/encabezado.png}}
  \end{picture}
  \begin{flushright}
    \vspace*{\fill}
    {\textcolor{titlecolor}{\LARGE\bfseries\MakeUppercase{\\[3cm]\@title}}}
    \vspace*{\fill}

    {\textcolor{titlecolor}{\Large\bfseries\scshape
      \begin{tabular}{@{}r@{}}
        José Basualto\\[0.2cm]
        Cristopher Dauros\\[0.2cm]
        Carlos Ruz
      \end{tabular}}}\\[1cm]
    {\textcolor{titlecolor}{\bfseries\textsc{
      Segundo informe\\[0.2cm]
      OII464: Desarrollo de Diseños Híbridos para Optimización\\[0.2cm]
      \@date}}}
  \end{flushright}
  \makeatother
\end{titlepage}

\hypersetup{pageanchor=true}
\pagenumbering{arabic}

\chapter*{Resumen}
La programación de cirugías electivas exige coordinar las etapas de anestesia
y cirugía, considerando la disponibilidad compartida de salas y profesionales
para evitar que la espera entre procedimientos prolongue la jornada o bloquee
recursos. Para abordar este problema, se propone integrar búsqueda tabú con
un bandido multibrazo dinámico (D-MAB), que aprende durante la ejecución
qué movimientos conviene utilizar para generar nuevos vecinos. El controlador
combina UCB1, encargado de equilibrar exploración y aprovechamiento de las
recompensas observadas, con Page--Hinkley, que detecta caídas de recompensa
y reinicia las estimaciones sin interrumpir la búsqueda tabú. La evaluación
contrastará esta metodología híbrida con Tabu de selección uniforme,
manteniendo iguales la inicialización, los movimientos disponibles y el
presupuesto de evaluaciones. Se utilizarán doce instancias sintéticas,
con treinta repeticiones por método e instancia, para analizar el objetivo
ponderado, el makespan, las esperas, el tiempo y la robustez,
respetando las restricciones de recursos y espera máxima. El diseño
permitirá estudiar si el controlador completo aporta mejoras reproducibles
y prácticamente relevantes, junto con su costo computacional, sin
presuponer una superioridad frente a la metaheurística base.

\noindent\textbf{Palabras clave:} programación de cirugías; búsqueda tabú;
selección adaptativa de operadores; aprendizaje en línea; bloqueo.

\clearpage
\section{Introducción}
Organizar una jornada quirúrgica supone distribuir una capacidad limitada
entre pacientes cuyas intervenciones tienen distintas duraciones y
requerimientos. La disponibilidad de quirófanos y personal especializado,
así como los tiempos de preparación y limpieza, condiciona tanto el orden
de atención como el término de la jornada. Cuando la anestesia y la cirugía se realizan en
etapas distintas, la espera por la segunda sala puede retener la primera
y su profesional, prolongando la jornada y dificultando el uso de los
recursos. En el problema considerado, se busca reducir el makespan,
entendido como el tiempo de término de la última operación, y el bloqueo,
respetando la precedencia, la espera máxima y la disponibilidad
\parencite{lara2026}. A diferencia de la suma de duraciones fijas, el
makespan depende de cómo se ordenan los pacientes y se asignan los recursos.

La literatura distingue roles del aprendizaje: \textcite{akbarzadeh2024}
predicen factibilidad en secuencias quirúrgicas multifase;
\textcite{niroumandrad2024} orientan Tabu con clasificadores para programar
médicos. En selección de operadores, \textcite{yu2024} emplean Q-learning
en programación y reprogramación quirúrgica, y el preprint de
\textcite{dolatkhah2026} usa D-MAB dentro de un algoritmo genético para
generar columnas iniciales. \textcite{yin2026} aportan Tabu iterado
multioperador entre hospitales, mientras \textcite{dacosta2008} constituyen
el fundamento histórico de D-MAB. Difieren en representación, restricciones
y evaluación; queda por establecer el aporte de D-MAB al Tabu del caso
docente de dos etapas con espera estricta, reparación y controles de cómputo.

El objetivo de este proyecto es estudiar si un controlador aprendido mejora
la calidad, el tiempo o la robustez de Tabu Search frente a la selección
uniforme, integrando una metaheurística (MH) con aprendizaje automático (ML).
D-MAB orientará la generación de vecinos durante la \emph{evolución} de la
trayectoria \parencite{vega_taxonomia}, mediante UCB1 y el detector
Page--Hinkley, mientras Tabu conservará la responsabilidad de aceptar
soluciones y mantener su memoria. La comparación se realizará en instancias
sintéticas del repositorio oficial \parencite{saicooh2026}, con las mismas
restricciones, soluciones iniciales y presupuestos, para evaluar el aporte
del controlador completo sin atribuir diferencias a un mayor esfuerzo de
búsqueda.

\clearpage
\section{Propuesta metodológica}
\subsection{Problema, representación y factibilidad}
La agenda será estática, con duraciones conocidas. La solución fija orden
y salas independientes por etapa; el planificador asigna profesionales
elegibles disponibles. La espera y el fin de jornada son
\begin{equation}
 W_j=\max\{0,T_j^C-F_j^A\},\qquad 0\leq W_j\leq L_j,
 \qquad C_{\max}=\max_{j\in J}F_j^C.
 \label{eq:restricciones}
\end{equation}
\noindent $J$ es el conjunto de pacientes y $j$ un paciente; $T_j^C$ es el
inicio de cirugía; $F_j^A$ y $F_j^C$ son los fines de cada etapa, incluida
limpieza; $W_j$ es la espera; $L_j$ su límite estricto, y $C_{\max}$ el
makespan. Las reservas siguen el modelo docente; la primera sala permanece
ocupada hasta el traslado si las salas difieren.

Se minimizará el siguiente criterio, sin penalización por desequilibrio de salas:
\begin{equation}
 f(S)=C_{\max}+10^{-6}\sum_o t_o^{\mathrm{ini}}
        +0{,}5\sum_{j\in J}W_j+1{,}4\max_{j\in J}W_j.
 \label{eq:objetivo}
\end{equation}
\noindent $S$ es un calendario y $f(S)$ su objetivo; $o$ recorre las
operaciones y $t_o^{\mathrm{ini}}$ es su inicio de reserva. $C_{\max}$, $J$
y $W_j$ se definieron en \eqref{eq:restricciones}. El coeficiente $10^{-6}$
desempata; $0{,}5$ y $1{,}4$ ponderan espera total y máxima. Se reportarán
estas métricas por separado: una reducción de $f$ no equivale necesariamente
a la misma reducción del makespan.

Las claves aleatorias representan orden y salas. El decodificador pospone
pacientes no programables y repara asignaciones. La reparación puede alterar
el movimiento y restringir calendarios alcanzables; se registrará su efecto
sobre la atribución de recompensas.

\subsection{Componente MH: búsqueda tabú}
Tabu parte de la mejor de treinta soluciones: FIFO (orden original),
SPT y LPT (duración total creciente y decreciente) y veintisiete aleatorias.
Propone intercambio, inserción y cambio de sala de anestesia,
cirugía o ambas. Explora lotes de quince candidatos, priorizando pacientes
con espera. Recuerda siete soluciones, con aspiración por mejora global;
si ninguna es admisible, usa la mejor del lote. Mantiene un archivo de
treinta soluciones, puede aceptar empeoramientos locales y conserva el
mejor resultado global.

\subsection{Componente ML: selección adaptativa en la evolución}
D-MAB elige el movimiento de cada candidato; cada operador es un brazo.
UCB1 aprende de las recompensas y combina su media con la exploración
de brazos menos probados. Page--Hinkley complementa esa decisión: detecta
caídas persistentes de recompensa y reinicia las estimaciones del bandido,
conservando el estado de Tabu \parencite{dacosta2008}.
La recompensa es la mejora positiva del candidato reparado frente al estado
local del lote, normalizada en $[0,1]$; los duplicados reciben cero sin
consumir evaluación. Se aprende durante la corrida, sin entrenamiento previo;
el planificador asegura factibilidad y Tabu selecciona el vecino aceptado.

\clearpage
\subsection{Orquestación del modelo híbrido}
Fuera de línea se prepararán y validarán las instancias y se fijarán los
parámetros con datos de desarrollo separados de las semillas de evaluación.
En línea, una única trayectoria Tabu generará lotes de candidatos con el
tipo de movimiento elegido por D-MAB. Después de reparar y evaluar cada
propuesta, su recompensa actualizará el controlador; Tabu elegirá el vecino
admisible y actualizará la memoria y la mejor solución global. Cada
ejecución terminará al alcanzar 6.060 evaluaciones, incluidas las treinta
iniciales, y devolverá su mejor calendario factible revalidado.

Se mantendrá Page--Hinkley como parte del controlador completo, porque la
utilidad de un movimiento puede cambiar durante la búsqueda. UCB1 realiza
el aprendizaje y la selección; Page--Hinkley es una prueba estadística
que indica cuándo renovar sus estimaciones. El detector no identifica
directamente un óptimo local ni garantiza que cada propuesta mejore.
Un candidato que empeora recibe recompensa cero, aunque Tabu puede aceptar
retrocesos locales para explorar sin perder su mejor solución.

La referencia será otra ejecución independiente del mismo Tabu, con
selección uniforme de los cinco movimientos y 6.060 evaluaciones.
Ambas metodologías compartirán el inicio de cada pareja instancia--semilla,
pero conservarán estados propios y se evaluarán por separado.
La comparación estudiará D-MAB como conjunto; los eventos del detector
se describirán sin atribuirle un efecto causal aislado.

\subsection{Datos, controles y plan de evaluación}
Se trabajará con 12 instancias obtenidas del repositorio oficial del problema
a estudiar, correspondientes a tres réplicas de 15, 20, 25 y 30 cirugías.
Se comprobarán esquema, hashes, elegibilidad y unidades temporales. El tamaño,
los recursos, las duraciones y la espera permitida son metacaracterísticas descriptivas;
no son entradas contextuales de D-MAB. Las nuevas semillas no sustituyen una
evaluación en instancias reservadas; la ampliación a escenarios de mayor
congestión y tamaños de 50, 100 y 150 cirugías se especificará como una fase
posterior, declarando su procedencia sintética.

El control uniforme sustituye la selección aprendida y permite evaluar
el aporte del controlador completo. FIFO, SPT y LPT se conservarán como
referencias constructivas e integrarán la inicialización común. La sección
\ref{sec:experimental} define las dos configuraciones, sus siglas y el
protocolo de comparación, con calidad, factibilidad, tiempo y robustez
medidos por separado. Se controlará especialmente la reparación, que puede
alterar el movimiento al que se atribuye la recompensa.

Las instancias no contienen identificadores clínicos. Se conservarán
calendarios finales, semillas, versiones y tiempos para permitir una
revisión reproducible y analizar el costo de la adaptación. La propuesta
se limita a una agenda estática con duraciones conocidas; urgencias,
cancelaciones y replanificación durante la jornada requieren una extensión
del modelo.

\clearpage
\section{Configuración experimental}
\label{sec:experimental}
El protocolo evaluará el aporte de D-MAB y de la cartera con controles de
igual presupuesto, siguiendo la comparación de algoritmos estocásticos
revisada en clases \parencite{vega_comparacion}. Se fijan los siguientes
criterios antes de ejecutar la campaña.

\subsection{Instancias y alcance}
Se utilizarán las 12 instancias sintéticas del repositorio oficial:
\texttt{HOSP-STD-N-01}, \texttt{-02} y \texttt{-03} para cada tamaño
de la Tabla~\ref{tab:instancias}. Se validarán esquema, elegibilidad, unidades
en minutos y huellas SHA-256. Se conservarán las restricciones y la función
objetivo de la propuesta, incluida la limpieza.

\begin{table}[ht]
 \centering
 \caption{Catálogo inicial y recursos disponibles.}
 \label{tab:instancias}
 \begin{tabular}{@{}rrrrrr@{}}
  \toprule
  Cirugías & Casos & Salas & Anestesistas & Cirujanos & Espera límite (min)\\
  \midrule
  15 & 3 & 6 & 5--6 & 7--10 & 18--55\\
  20 & 3 & 7 & 6--7 & 9--12 & 19--61\\
  25 & 3 & 8 & 7--8 & 12--15 & 21--65\\
  30 & 3 & 9 & 8--9 & 15--18 & 23--69\\
  \bottomrule
 \end{tabular}
\end{table}
Los intervalos reflejan variación entre casos y, para el límite de espera,
también entre pacientes. Este catálogo ya es conocido por el proyecto:
las nuevas semillas medirán variabilidad de búsqueda dentro de él.
La ampliación a 50, 100 y 150 cirugías se registrará como una fase separada,
con reglas de generación y admisión fijadas antes de comparar métodos;
esas instancias se reservarán de la calibración.

\subsection{Métodos y controles}
Se mantendrán dos bloques (Tabla~\ref{tab:metodos}). En A se aislarán la
selección aprendida y el detector; en B se evaluará la cartera frente
a prolongar Tabu o repetirlo. FIFO, SPT y LPT se calcularán una vez por
instancia como referencias deterministas, además de integrar la inicialización.

\begin{table}[ht]
 \centering
 \caption{Configuraciones de la campaña principal.}
 \label{tab:metodos}
 \begin{tabular}{@{}lp{0.65\textwidth}r@{}}
  \toprule
  Bloque & Método e identificador & Evaluaciones\\
  \midrule
  A & U: Tabu con selección uniforme & 3.030\\
  A & M: Tabu con D-MAB & 3.030\\
  A & Q: Tabu con UCB1 sin Page--Hinkley & 3.030\\
  B & T: Tabu uniforme prolongado & 6.060\\
  B & R: dos Tabu uniformes independientes & $3.030+3.030$\\
  B & H: cartera Tabu uniforme + Tabu con D-MAB & $3.030+3.030$\\
  \bottomrule
 \end{tabular}
\end{table}
La comparación estadística se efectuará dentro de cada bloque. La conservación
del camino uniforme de H no implica superioridad frente a T o R.

\clearpage
\subsection{Repeticiones, parámetros y parada}
Cada configuración tendrá 30 repeticiones por instancia, con semillas
80--109: se obtendrán 2.160 registros método--instancia--semilla.
U, M, Q y T compartirán las treinta soluciones iniciales de cada semilla;
H conservará el mismo inicio para sus dos caminos. R usará las semillas
$s$ y $s+1.000.000$, donde $s$ es la semilla principal.
El pareamiento se justificará por este diseño y por la instancia compartida;
los sorteos posteriores pueden divergir.

Se fijan 30 soluciones iniciales, 15 candidatos por lote y memoria de siete
soluciones. En M y Q se usará exploración $C=0{,}1$; M tendrá
margen $\delta=0{,}01$ y umbral $\lambda=1$ en Page--Hinkley.
No se ajustarán estos valores con las semillas de evaluación.
Cualquier calibración posterior utilizará datos de desarrollo y una versión
distinta del protocolo.

El presupuesto contará las valoraciones de candidatos, incluidas las
treinta iniciales de cada camino. Los duplicados sin reevaluación no gastarán
otra valoración; sus intentos y la reparación se registrarán.
La revalidación final se contabilizará aparte. Toda solución devuelta
deberá cumplir las restricciones estrictas; los fallos se conservarán en el
registro y nunca se sustituirán silenciosamente por otra semilla.

Como análisis secundario, B se comparará con 10 segundos de búsqueda:
T dispondrá de 10 s, y R y H de dos cuotas secuenciales de 5 s,
con inicialización incluida y sin transferir sobrantes.
Un plazo global de 10 s limitará ambas cuotas: el sobrepaso del primer
camino se descontará del segundo. El reloj monotónico se comprobará antes
de generar cada candidato; se registrarán el sobrepaso atómico final y las
inicializaciones incompletas. Este criterio reemplazará al de evaluaciones.
Se medirán por separado búsqueda y ejecución total, incluida la validación.

\subsection{Métricas y relevancia práctica}
La métrica primaria será $f(S)$ de \eqref{eq:objetivo}; makespan, espera
total y máxima, factibilidad, tiempo, número de evaluaciones y convergencia
se reportarán por separado. Se describirán mediana e intervalo intercuartílico,
junto con frecuencias de movimientos, reparación y reinicios del controlador.
Las curvas usarán puntos comunes del presupuesto: 10, 50 y 100\,\%.

Para comparar escalas se calculará
\begin{equation}
 g_i=100\frac{m_{c,i}-m_{a,i}}{m_{c,i}},
 \label{eq:ganancia}
\end{equation}
\noindent donde $i$ es una instancia; $m_{c,i}$ y $m_{a,i}$ son las medianas
de las treinta corridas del control y la alternativa, respectivamente,
y $g_i$ es su mejora relativa porcentual; $g_i>0$ favorece a la alternativa.
El efecto mínimo relevante se fija en \textbf{1\,\% del objetivo ponderado},
sin equipararlo a un porcentaje de minutos. La tolerancia numérica
$10^{-9}$ se aplicará a diferencias para identificar empates, conservando
los objetivos crudos. El umbral práctico y la significación se interpretarán
por separado, junto con sus intervalos de confianza.

\clearpage
\subsection{Hipótesis e inferencia estadística}
En la campaña por evaluaciones se fijan cuatro contrastes globales bilaterales:
M frente a U, H frente a T,
M frente a Q y H frente a R. El nivel será $\alpha=0{,}05$,
con corrección de Holm conjunta para las cuatro comparaciones.
La hipótesis nula de Wilcoxon será simetría de las diferencias
alrededor de cero, frente a la alternativa bilateral.
Se usará una mediana por instancia y método: \textbf{doce bloques globales},
sin tratar las 360 corridas de un método como instancias independientes.

Para inferencia se dividirá la diferencia entre ambas medianas por su máximo;
esta normalización conserva la magnitud al intercambiar métodos.
Wilcoxon usará las doce diferencias normalizadas, con permutaciones de signos
y tratamiento explícito de ceros y empates
\parencite{scipy_wilcoxon,scipy_permutation}.
Se fijará la prueba de signos como sensibilidad ante asimetría;
se reportarán ambas sin elegir la que dé menor valor $p$.
Por instancia se aplicará la misma normalización a las treinta parejas;
sus cuatro contrastes por doce instancias constituirán una familia exploratoria
de 48 pruebas, corregida con Holm.
La secundaria temporal usará las mismas semillas; H frente a T y H frente a R
formarán una familia exploratoria separada con Holm.

Como análisis contextual se aplicará Friedman a los tres métodos de cada bloque.
Por el tamaño de la muestra se usarán 100.000 permutaciones de etiquetas
dentro de cada instancia \parencite{scipy_friedman,scipy_permutation}.
Se corregirán con Holm los dos contrastes ómnibus.
Solo ante rechazo se interpretarán las tres comparaciones post hoc del bloque,
con Holm; estas no sustituirán los cuatro contrastes globales prefijados.

Se estimará la mediana de las mejoras relativas y su IC95\,\% con 5.000
remuestreos: instancias y, dentro de ellas, semillas emparejadas, conservando
juntos los métodos. Se añadirá $A_{12}$ por instancia como medida descriptiva
de superioridad marginal: más de 0,5 favorecerá a la alternativa al minimizar.
Las semillas estadísticas serán independientes y se fijarán en el protocolo.
No rechazar una hipótesis no demostrará equivalencia; treinta repeticiones
tampoco garantizan potencia suficiente para detectar una mejora del 1\,\%.

\subsection{Entorno, registros y reproducibilidad}
La referencia será macOS/arm64, Python 3.12.14, NumPy 2.3.5 y SciPy 1.18.1.
Se registrarán CPU, memoria y versiones efectivas, ejecutando secuencialmente
en la misma máquina y sin otras tareas de cálculo concurrentes.
Cada corrida guardará instancia, método, semilla, parámetros, evaluaciones,
tiempos, calidad, factibilidad y calendario final, además de trazas y eventos
del controlador.

El manifiesto incluirá huellas de datos y de todos los archivos ejecutados,
junto con la revisión del código. Se verificarán cobertura, duplicados y
calendarios antes del análisis. Los casos fallidos se describirán aparte,
con su tasa de fallos; una instancia sin resúmenes comparables no entrará
silenciosamente en una prueba.
Los contrastes exigirán treinta corridas válidas por método e instancia;
si falta alguna, se declarará incompleta la comparación correspondiente.
Se conservarán los datos sintéticos y las versiones del protocolo;
cualquier enmienda se declarará antes de su campaña.


\clearpage
\section{Conclusión}
El diseño propuesto integra aprendizaje en la evolución de Tabu Search:
UCB1 selecciona movimientos a partir de su recompensa y Page--Hinkley
renueva las estimaciones cuando detecta una caída, mientras la
metaheurística mantiene la memoria, la aceptación de vecinos y la mejor
solución. Esta distribución de responsabilidades permite estudiar un
controlador adaptativo completo para programar anestesia y cirugía,
respetando disponibilidad, limpieza y espera máxima, con el propósito
de reducir el objetivo ponderado y examinar sus efectos sobre el makespan
y el bloqueo.

La comparación con Tabu uniforme mantendrá iguales el inicio, los
movimientos y el presupuesto, acompañando las diferencias de calidad
con medidas de incertidumbre, tiempo y robustez. El protocolo fija una
mejora relevante del 1\,\% del objetivo y registra la reparación y los
reinicios para interpretar el comportamiento observado, sin anticipar
superioridad ni atribuir un beneficio aislado al detector. La ampliación
posterior a instancias reservadas y de mayor congestión permitirá
estudiar el alcance de las conclusiones y el costo de aumentar la escala.

\clearpage
\printbibliography[title={Referencias}]
\end{document}
